AI tools were used heavily in this project, and this section states where.

\begin{itemize}
\item \textbf{My own work.} I chose the topic, wrote the first serial prototype
of attention and online softmax (\code{src/attention.cpp}, \code{src/main.cpp}),
decided the scope and priorities of the study, arranged the cluster accounts and
access, and reviewed the results and this report.
\item \textbf{Claude (Anthropic), in the chat interface and as the Claude Code
agent.} Early on it was used for brainstorming the topic, planning experiments and
reviewing my prototype. In the final phase the Claude Code agent wrote most of the
production code under my direction: the tuned serial kernel, the no-communication
control, the CUDA kernels and the MPI+CUDA driver, the microbenchmark and analysis
programs, the extraction script, the cluster job scripts and the plotting scripts.
It also ran the experiments on Vanda through my account (job scripts, submission,
monitoring and retrieval of results), diagnosed the two measurement problems of
\cref{sec:pitfalls} from the data, produced the figures and tables, and drafted the
text of this report.
\item \textbf{Codex (OpenAI).} It reviewed the project plan, wrote the first version
of the flat-tensor compute core, the MPI variants, the benchmark driver and the
tests, and performed a code review that found a numerical error in an early version
of the vectorised exponential (\cref{sec:validation}). Near the end it reviewed the
code, the retrieved data and a draft of this report; that review found two edge-case
errors in the analysis program and several claims that were stronger than the
evidence, and Claude Code then made the corresponding changes.
\item \textbf{Later review rounds.} After that, the report was checked several
more times against the code and the result files, in some rounds by me and in
others by a separate Claude session. These rounds led to the statements about the balance of threads under a causal
mask and about the query tiles of threaded ranks (\cref{sec:xdesign}), to the
set-up cost of the one-sided variants (\cref{sec:setup}), to the paragraph on
repeatability (\cref{sec:validation}), to the result sentence in the figure
captions and to several corrections of wording; Claude Code made these changes
as well.
\end{itemize}

\noindent Every measured value in the report comes from a job that actually ran on
the cluster (\cref{tab:provenance}) and is computed from the result files by a
script rather than typed. The correctness of the AI-written code does not rest on
trust: each implementation is compared with a double-precision reference inside
every timed scaling run (the runs that measure memory, hardware counters or the
network carry no comparison), the unit tests ran on the cluster on the submitted
source, and the real-data path is additionally compared with PyTorch's own
output.
